% !TEX program = xelatex
% 공통수학2 · Ⅱ. 집합과 명제 · 중단원(2.명제)·대단원(Ⅱ) 마무리
\documentclass[11pt,a4paper]{article}
\usepackage{kotex}
\usepackage{iftex}
\ifPDFTeX\else
  \setmainhangulfont{Noto Serif CJK KR}
  \setsanshangulfont{Noto Sans CJK KR}
\fi
\usepackage[margin=2.1cm]{geometry}
\usepackage{parskip}
\usepackage{enumitem}
\usepackage{array}
\usepackage{tabularx}
\usepackage{booktabs}
\usepackage{xcolor}
\usepackage{amssymb}
\usepackage{tikz}
\usepackage[most]{tcolorbox}
\usepackage{fancyhdr}
\usepackage{titlesec}

\definecolor{goldc}{HTML}{B07C22}
\definecolor{goldstrong}{HTML}{8A5F16}
\definecolor{indigoc}{HTML}{2B3B57}
\definecolor{indigosoft}{HTML}{6E82A6}
\definecolor{linec}{HTML}{D6DACB}
\definecolor{surface2}{HTML}{E4E8DC}
\definecolor{notebg}{HTML}{FBF3E3}
\definecolor{inksoft}{HTML}{52604F}

\titleformat{\section}{\normalfont\sffamily\bfseries\large\color{indigoc}}{\thesection}{0.6em}{}
\titlespacing{\section}{0pt}{16pt}{8pt}
\newtcolorbox{ideabox}{enhanced, breakable, colback=white, colframe=goldc, boxrule=0.5pt, leftrule=3pt, arc=2pt, left=10pt, right=10pt, top=8pt, bottom=8pt}
\newtcolorbox{calloutbox}{enhanced, breakable, colback=white, colframe=indigosoft, boxrule=0.5pt, leftrule=3pt, arc=2pt, left=10pt, right=10pt, top=7pt, bottom=7pt}
\newtcolorbox{notebox}{enhanced, breakable, colback=notebg, colframe=goldc, boxrule=0.5pt, leftrule=3pt, arc=2pt, left=10pt, right=10pt, top=7pt, bottom=7pt}
\newtcolorbox{answerbox}{enhanced, breakable, colback=white, colframe=linec, boxrule=0.5pt, arc=2pt, left=10pt, right=10pt, top=7pt, bottom=7pt}
\newcommand{\lessonbanner}[3]{%
  \begin{tcolorbox}[enhanced, colback=indigoc, colframe=indigoc, boxrule=0pt, arc=0pt,
    left=14pt, right=14pt, top=14pt, bottom=10pt, borderline south={2.4pt}{0pt}{goldc}, fontupper=\color{white}]
    {\sffamily\footnotesize\color{white!85!black} #1}\\[4pt]{\sffamily\bfseries\Large #2}\\[3pt]{\sffamily\small\color{white!88!black} #3}
  \end{tcolorbox}}
\pagestyle{fancy}\fancyhf{}\renewcommand{\headrulewidth}{0pt}
\fancyfoot[L]{\footnotesize\color{inksoft} 공통수학2 · 2026학년도 2학기 · 지평선고등학교}
\fancyfoot[R]{\footnotesize\color{inksoft} 교사용 지도안 -- 배포 금지}

\begin{document}
\lessonbanner{공통수학2 · Ⅱ. 집합과 명제}{중단원(2. 명제) \textperiodcentered\ 대단원(Ⅱ) 마무리 --- 32차시}{명제 전체 종합 \textperiodcentered\ 교사용 지도안 (2026학년도 2학기)}
\vspace{10pt}

\noindent\begin{tabularx}{\linewidth}{@{}>{\bfseries\color{indigoc}}p{2.6cm}X@{}}
\toprule
대상 & 고1 · 2개 반 · 반당 13명 \\
차시 & 50분 (도입 10 · 전개 30 · 정리 10) \\
성취기준 & 명제와 조건, 명제 사이의 관계, 명제의 증명을 통합적으로 이해하고 적용할 수 있다. \\
선행 학습 & 26$\sim$31차시 --- `2. 명제' 전체 (명제와 조건 / 명제 사이의 관계 / 명제의 증명) \\
\bottomrule
\end{tabularx}

\section{학습 목표 \& Big Idea}
\begin{ideabox}
\textbf{\color{goldstrong}영속적 이해 (Big Idea)}\\
`집합과 명제'라는 대단원 전체를 관통하는 하나의 아이디어는, 참·거짓이라는 논리적 판단이 결국 진리집합 사이의 \textbf{포함 관계}로 환원된다는 것이다. 명제의 참·거짓, 역과 대우, 충분·필요조건, 증명법은 모두 이 하나의 원리에서 갈라져 나온다.
\vspace{4pt}
\small\color{inksoft} 핵심개념: \textbf{포함 관계} \quad\textbullet\quad 핵심개념: \textbf{논리적 동치} \quad\textbullet\quad 일반화: \textbf{집합의 언어가 명제의 언어를 뒷받침한다}
\end{ideabox}
\begin{calloutbox}
\textbf{차시 학습 목표}
\begin{itemize}[leftmargin=1.4em, itemsep=2pt, topsep=4pt]
  \item 명제와 조건, 진리집합, 역과 대우, 충분·필요조건, 증명법을 종합 문제를 통해 점검하고 정리할 수 있다.
  \item 대단원 Ⅱ `집합과 명제'에서 배운 개념들이 어떻게 서로 연결되는지 설명할 수 있다.
\end{itemize}
\end{calloutbox}

\section{질문의 연결}
\begin{tabularx}{\linewidth}{@{}>{\bfseries\color{indigoc}\small}p{2.4cm}X@{}}
사실적 질문 & `1. 집합'과 `2. 명제'에서 배운 개념들을 각각 한 문장으로 요약하면 무엇일까? \\[6pt]
개념적 질문 & 집합에서 배운 포함 관계($\subset$)는 명제의 참·거짓, 충분·필요조건과 어떻게 하나로 연결될까? \\[6pt]
논쟁적 질문 & 수학에서 `증명된 것'과 일상에서 `설득력 있는 것'은 같은 것일까, 다른 것일까? \\
\end{tabularx}

\section{수업 흐름}
\begin{tabularx}{\linewidth}{@{}>{\centering\arraybackslash}p{1.6cm}X>{\raggedright\arraybackslash\small}p{3.1cm}@{}}
\toprule
\textbf{단계} & \textbf{활동 \& 발문} & \textbf{ATL / VTR} \\
\midrule
\textcolor{goldstrong}{\textbf{도입}}\newline\footnotesize 10분 &
\textbf{마음공부 (2분)} --- 대단원을 마무리하는 유무념 대조.\newline
\textbf{개념 지도 그리기 (8분)} --- `집합과 명제' 대단원의 하위 개념(집합의 연산, 명제와 조건, 역과 대우, 충분·필요조건, 증명법)을 화이트보드에 자유롭게 배치하며 서로 어떻게 연결되는지 화살표로 표시. &
ATL: 사고(구조화)\newline VTR: Connect-Extend-Challenge \\
\addlinespace
\textcolor{goldstrong}{\textbf{전개}}\newline\footnotesize 30분 &
\textbf{종합 문제 풀이 (30분)} --- 문제 1(명제 판별) $\to$ 문제 2(진리집합) $\to$ 문제 3(역·대우와 참·거짓) $\to$ 문제 4(충분·필요·필요충분조건)를 모둠별로 순환하며 해결. 각 문제를 풀 때마다 `이 문제가 어느 차시의 개념과 연결되는지' 짚어 준다. &
ATT: 촉진자 역할\newline ATL: 협업(모둠 순환)\newline VTR: 개념 연결 확인 \\
\addlinespace
\textcolor{goldstrong}{\textbf{정리}}\newline\footnotesize 10분 &
\textbf{개념 연결 질문 (5분)} --- ``대단원 Ⅱ에서 가장 헷갈렸던 개념은 무엇이었고, 오늘 다시 보니 어떻게 정리되었나?''\newline
\textbf{성찰 (5분)} --- 감사 일기 1문장 + 대단원 Ⅱ `집합과 명제' 총마무리 + 다음 차시 예고(대단원 Ⅲ `함수와 그래프' 시작, `함수'). &
마음공부 · 감사 일기 \\
\bottomrule
\end{tabularx}

\begin{calloutbox}
\textbf{지도상의 유의점} --- 이 차시는 새로운 개념을 배우는 시간이 아니라 이미 배운 개념들 사이의 \textbf{연결}을 확인하는 시간임을 분명히 한다. 학생이 특정 문제를 틀렸을 때는 정답만 알려 주지 않고 `이 문제는 몇 차시의 어떤 개념과 연결되는지'를 함께 되짚어 스스로 근거를 찾도록 돕는다.
\end{calloutbox}

\section{형성평가 \& 답안}
\begin{minipage}[t]{0.48\linewidth}
\begin{answerbox}
\textbf{문제 1} \quad 명제 판별\\
\small ⑴ $\sqrt4$는 유리수이다. ⑵ $x$는 100보다 큰 수이다. ⑶ 26은 3의 배수이다. ⑷ 우리나라는 아름다운 나라이다.\\[4pt]
\textbf{\color{goldstrong}답} \; 명제: ⑴(참, $\sqrt4=2$), ⑶(거짓) \quad ⑵는 조건(변수 포함), ⑷는 주관적 판단이라 명제 아님
\end{answerbox}
\end{minipage}\hfill
\begin{minipage}[t]{0.48\linewidth}
\begin{answerbox}
\textbf{문제 2} \quad $U=\{x\mid x$는 7 이하의 자연수$\}$, $p:x^2-6x+8=0$, $q:x^2-6x+9>0$의 진리집합\\[4pt]
\textbf{\color{goldstrong}답} \; ⑴ $P=\{2,4\}$ \; ⑵ $Q=\{1,2,4,5,6,7\}$ \; ⑶ $P^C=\{1,3,5,6,7\}$ \; ⑷ $Q^C=\{3\}$
\end{answerbox}
\end{minipage}

\vspace{6pt}
\begin{minipage}[t]{0.48\linewidth}
\begin{answerbox}
\textbf{문제 3} \quad 역과 대우, 참·거짓\\
\small ⑴ $x$가 자연수이면 $x$는 유리수이다.\\
⑵ $x(x+1)=0$이면 $x+1=0$이다.\\[4pt]
\textbf{\color{goldstrong}답} \; ⑴ 참. 역(거짓, $x=\frac12$가 반례) / 대우(참) \quad ⑵ 거짓($x=0$이 반례). 역(참) / 대우(거짓)
\end{answerbox}
\end{minipage}\hfill
\begin{minipage}[t]{0.48\linewidth}
\begin{answerbox}
\textbf{문제 4} \quad 어떤 조건인지 판별\\
\small ⑴ $p:x\le1$, $q:-3\le x\le0$\\
⑵ $p:a=b$, $q:a^2=b^2$\\
⑶ $p:|x|\le1$, $q:-1\le x\le1$\\[4pt]
\textbf{\color{goldstrong}답} \; ⑴ 필요조건 \; ⑵ 충분조건 \; ⑶ 필요충분조건
\end{answerbox}
\end{minipage}

\section{성취수준 (대단원 Ⅱ 종합)}
\begin{tabularx}{\linewidth}{@{}>{\bfseries\color{goldstrong}}p{0.9cm}X@{}}
\toprule
A & 집합과 명제의 모든 개념을 통합적으로 이해하고, 근거를 들어 스스로 설명하며 다양한 문제에 적용할 수 있다. \\
B & 집합과 명제의 주요 개념을 이해하고 대부분의 문제를 스스로 해결할 수 있다. \\
C & 집합과 명제의 기본 개념을 알고, 안내에 따라 문제를 해결할 수 있다. \\
D & 집합과 명제의 일부 개념을 알고, 안내된 절차에 따라 간단한 문제를 해결할 수 있다. \\
E & 집합과 명제의 기본 용어를 안다. \\
\bottomrule
\end{tabularx}

\section{다음 차시}
\begin{calloutbox}
\textbf{대단원 Ⅲ. 함수와 그래프 --- 1. 함수 (33차시).} 대단원 Ⅱ `집합과 명제'가 여기서 마무리된다. 다음 시간부터는 함수의 뜻과 그래프를 다룬다.
\end{calloutbox}
\end{document}
